\chapter{Los conjuntos y las variantes}\label{ch:parametros}

\begin{roadmap}
El capítulo~\ref{ch:keygen} mostró cómo nace una clave; los
capítulos~\ref{ch:encapdecap} y~\ref{ch:labkem}, cómo se encapsula y
desencapsula. Falta el catálogo: qué conjuntos concretos define Classic
McEliece, qué significan los sufijos \texttt{f}, \texttt{pc} y
\texttt{pcf}, cómo se guardan claves y criptogramas byte a byte, y qué
seguridad se les atribuye. Al terminar, el lector debería poder elegir
un conjunto con criterio y saber cuántos bytes maneja en cada paso.
\end{roadmap}

\section{Los cinco tamaños}\label{sec:param-cinco}

Las tablas~\ref{tab:keygen-pol} y~\ref{tab:keygen-tamanos} fijaron los
cinco pares $(m,n,t)$ y el desglose interno de las claves. La
tabla~\ref{tab:param-base} añade las dos cifras que faltaban para
decidir: el criptograma y la categoría de seguridad.

\begin{table}[htbp]
\centering
\small
\setlength{\tabcolsep}{7pt}
\begin{tabular}{lcrrrc}
\toprule
Conjunto & categoría & $ek$ (bytes) & $ct$ (bytes) & $sk$ (bytes) & ¿ISO?\\
\midrule
\texttt{mceliece348864} & 1 & \num{261120}  & 96  & \num{6492}  & no\\
\texttt{mceliece460896} & 3 & \num{524160}  & 156 & \num{13608} & sí\\
\texttt{mceliece6688128} & 5 & \num{1044992} & 208 & \num{13932} & sí\\
\texttt{mceliece6960119} & 5 & \num{1047319} & 194 & \num{13948} & sí\\
\texttt{mceliece8192128} & 5 & \num{1357824} & 208 & \num{14120} & sí\\
\bottomrule
\end{tabular}
\caption{Los cinco conjuntos base, con
$ek=mt\lceil k/8\rceil$, $ct=\lceil mt/8\rceil$ y $sk$ como en la
tabla~\ref{tab:keygen-tamanos}. La columna «¿ISO?» anticipa la
sección siguiente: el estándar ISO incluye los cuatro tamaños
mayores, cada uno con cuatro variantes; \texttt{mceliece348864} queda
fuera.}\label{tab:param-base}
\end{table}

¿De dónde salen estos números, a primera vista caprichosos? Los
conjuntos se eligieron siempre \emph{maximizando la seguridad sujeta a
un límite de tamaño} \cite{cmrationale}:

\begin{itemize}[leftmargin=1.5em,itemsep=3pt]
  \item \texttt{mceliece6960119}, con $(n,t)=(6960,119)$, maximiza la
        seguridad para claves públicas que quepan en $2^{20}$ bytes.
        La elección de $n=6960$ viene de la optimización clásica de
        2008: para límites de $2^{16}$ a $2^{20}$ bytes, los códigos
        Goppa de longitudes $1744$, $2480$, $3408$, $4624$ y $6960$,
        con $t$ proporcionalmente creciente, maximizan la dificultad
        del ataque; la tasa óptima ronda $0{,}75$. La única diferencia
        con aquella elección es que Classic McEliece añade $119$
        errores en lugar de $121$: sin decodificación por listas no
        puede decodificar de forma única más allá de la capacidad del
        código, y añadir dos errores menos reduce marginalmente la
        seguridad.
  \item \texttt{mceliece6688128} también maximiza dentro de $2^{20}$
        bytes, pero exigiendo que $n$ y $t$ sean múltiplos de $32$;
        \texttt{mceliece460896} y \texttt{mceliece348864} hacen lo
        mismo dentro de $2^{19}$ y $2^{18}$ bytes.
  \item \texttt{mceliece8192128} es el conjunto «clásico»: $n$ y $t$
        potencias de dos, como en el McEliece original de 1978, a
        costa de pasarse un poco de $2^{20}$ bytes
        ($\num{1357824}$ bytes $=1{,}29$~MiB).
\end{itemize}

Maximizar bajo un límite es más robusto que apuntar a un nivel exacto:
las métricas de seguridad varían entre modelos, y un conjunto clavado
a un objetivo puede quedar por debajo en otro modelo distinto. Al
maximizar por tamaño, la incertidumbre cambia poco el orden.

\begin{observacion}{Múltiplos de $32$: por qué}{param-mult32}
Cuando $n$ y $t$ son múltiplos de $32$ ---y por tanto $mt$ y $k$
también lo son, porque $k=n-mt$ y $q=2^{13}$--- no hay \emph{bits de
relleno} en las codificaciones de claves y criptogramas, y la
distinción entre las dos formas de leer un criptograma
(sección~\ref{sec:param-narrowly}) desaparece. Es una pequeña ventaja
de interoperabilidad a cambio de un pequeño empeoramiento en la
relación tamaño-seguridad. Los conjuntos $3488/4608/6688$ y el
$8192$ la tienen; el $6960$, no.
\end{observacion}

\section{Los dieciséis conjuntos de ISO}\label{sec:param-iso}

ISO estandarizó Classic McEliece en junio de 2026 (enmienda 2 de
ISO/IEC 18033-2:2006 \cite{iso18033}), y el estándar incluye
\textbf{dieciséis conjuntos}: los cuatro tamaños de la
tabla~\ref{tab:param-base} (sin \texttt{mceliece348864}) cruzados con
cuatro variantes:

\begin{itemize}[leftmargin=1.5em,itemsep=2pt]
  \item el conjunto \textbf{base}, tal cual lo define el equipo;
  \item el sufijo \textbf{\texttt{f}}: forma semi-sistemática
        $(\mu,\nu)=(32{,}64)$ (sección~\ref{sec:param-f});
  \item el sufijo \textbf{\texttt{pc}}: confirmación de texto claro
        (sección~\ref{sec:param-pc});
  \item el sufijo \textbf{\texttt{pcf}}: ambas cosas a la vez.
\end{itemize}

La tabla~\ref{tab:param-iso} y la figura~\ref{fig:param-iso} muestran
el catálogo completo. Obsérvese lo que cada sufijo hace a los tamaños:
\texttt{f} no cambia nada; \texttt{pc} suma exactamente
$\lceil\ell/8\rceil=32$ bytes al criptograma.

\begin{table}[htbp]
\centering
\small
\setlength{\tabcolsep}{6pt}
\begin{tabular}{lcccccc}
\toprule
tamaño & $(m,n,t,k)$ & $ek$ & $sk$ & $ct$ \texttt{base},\texttt{f} & $ct$ \texttt{pc},\texttt{pcf} & cat.\\
\midrule
\texttt{460896}  & $(13,4608,96,3360)$  & \num{524160}  & \num{13608} & 156 & 188 & 3\\
\texttt{6688128} & $(13,6688,128,5024)$ & \num{1044992} & \num{13932} & 208 & 240 & 5\\
\texttt{6960119} & $(13,6960,119,5413)$ & \num{1047319} & \num{13948} & 194 & 226 & 5\\
\texttt{8192128} & $(13,8192,128,6528)$ & \num{1357824} & \num{14120} & 208 & 240 & 5\\
\bottomrule
\end{tabular}
\caption[Los dieciséis conjuntos de ISO, cada fila con sus cuatro variantes.]{Los dieciséis conjuntos de ISO: cada fila, con los sufijos
---, \texttt{f}, \texttt{pc}, \texttt{pcf}. El nombre completo es
\texttt{mceliece} seguido del número y el sufijo, por ejemplo
\texttt{mceliece6960119pcf}. Los tamaños de cada fila valen para las
cuatro variantes salvo el criptograma de \texttt{pc}/\texttt{pcf}.}
\label{tab:param-iso}
\end{table}

\begin{figure}[htbp]
\centering
\begin{tikzpicture}[
    chip/.style={draw=mcblue,fill=mcsky!45,rounded corners=3pt,align=center,
      font=\footnotesize,inner sep=4pt,minimum width=3.5cm,minimum height=1.05cm},
    pchip/.style={draw=mcorange,fill=mcorange2!55,rounded corners=3pt,align=center,
      font=\footnotesize,inner sep=4pt,minimum width=3.5cm,minimum height=1.05cm},
    head/.style={font=\small\bfseries}
  ]
  \node[head,color=mcblue] at (2.0,0.85) {base};
  \node[head,color=mcblue] at (6.0,0.85) {\texttt{f}};
  \node[head,color=mcorange] at (10.0,0.85) {\texttt{pc}};
  \node[head,color=mcorange] at (14.0,0.85) {\texttt{pcf}};
  \foreach \y/\nm/\ca/\cb in {0/mceliece460896/156/188,-1.4/mceliece6688128/208/240,-2.8/mceliece6960119/194/226,-4.2/mceliece8192128/208/240} {
    \node[chip] at (2.0,\y) {\texttt{\nm}\\[1pt] $ct$: \ca\ bytes};
    \node[chip] at (6.0,\y) {\texttt{\nm}\texttt{f}\\[1pt] $ct$: \ca\ bytes};
    \node[pchip] at (10.0,\y) {\texttt{\nm}\texttt{pc}\\[1pt] $ct$: \cb\ bytes};
    \node[pchip] at (14.0,\y) {\texttt{\nm}\texttt{pcf}\\[1pt] $ct$: \cb\ bytes};
  }
  \draw[decorate,decoration={brace,amplitude=5pt},thick,mcorange]
    (8.2,-4.95) -- (15.8,-4.95);
  \node[font=\scriptsize,color=mcorange,align=center] at (12.0,-5.55)
    {$+\lceil\ell/8\rceil=32$ bytes de confirmación};
  \draw[decorate,decoration={brace,mirror,amplitude=5pt},thick,mcblue]
    (4.2,-4.95) -- (7.8,-4.95);
  \node[font=\scriptsize,color=mcblue,align=center] at (6.0,-5.55)
    {mismos tamaños:\\ otra generación de clave};
\end{tikzpicture}
\caption{Los dieciséis conjuntos de ISO, como familia. Los cuatro
tamaños (filas) y las cuatro variantes (columnas). La confirmación de
texto claro (\texttt{pc}, \texttt{pcf}) añade $32$ bytes al
criptograma; la variante \texttt{f} no cambia ningún tamaño, solo la
manera de generar la clave (forma semi-sistemática).}
\label{fig:param-iso}
\end{figure}

¿Por qué falta \texttt{mceliece348864}? Por el modelo cuántico de ISO.
ISO pidió a los candidatos \textbf{128 bits de seguridad en su modelo
cuántico}, que asume la aceleración de raíz cuadrada de la búsqueda de
Grover; pedir $128$ bits en ese modelo es mucho más exigente que pedir
$128$ bits clásicos \cite{cmiso}. Con la raíz cuadrada aplicada a los
mejores ataques clásicos sin trucos de memoria:

\begin{center}
\small
\begin{tabular}{lcccc}
\toprule
conjunto & ataque clásico & mitad (modelo cuántico) & ¿$\ge2^{128}$?\\
\midrule
\texttt{mceliece348864} & $2^{151{,}4}$ & $2^{75{,}7}$ & no\\
\texttt{mceliece460896} & $2^{193{,}3}$ & $2^{96{,}7}$ & no\\
\texttt{mceliece6688128} & $2^{268{,}0}$ & $2^{134{,}0}$ & sí\\
\texttt{mceliece6960119} & $2^{268{,}4}$ & $2^{134{,}2}$ & sí\\
\texttt{mceliece8192128} & $2^{303{,}2}$ & $2^{151{,}6}$ & sí\\
\bottomrule
\end{tabular}
\end{center}

Los valores clásicos son los de Stern bajo memoria gratis
(tabla~\ref{tab:param-seguridad}), cuyos estimadores cuentan
operaciones vectoriales como una sola operación, de modo que son cotas
prudentes. Con esa vara, el borrador ISO que el equipo publicó en
2023 incluyó solo los tamaños $6688$, $6960$ y $8192$ ---todos
$\ge2^{134}$ cuánticos---. Pero la aceleración de Grover está, en la
práctica, limitada por la latencia del ataque, así que ISO decidió
estandarizar también \texttt{mceliece4*} (el conjunto $460896$) a
sabiendas de que no alcanza el listón cuántico
\cite{cmiso}. El equipo recomienda los tamaños
\texttt{mceliece6*} para seguridad a largo plazo.

\begin{historia}[Cómo llegó Classic McEliece a ISO]
En octubre de 2022, unas diapositivas de Microsoft revelaron que el
subcomité ISO/IEC JTC1/SC27/WG2 había solicitado comentarios sobre la
inclusión de FrodoKEM, Classic McEliece y CRYSTALS-Kyber en sus
estándares. El equipo de Classic McEliece preparó como borrador el
documento \texttt{iso-mceliece-20230419} y lo publicó para recibir
comentarios. Siguiendo la regla ISO de expresar los requisitos «en
términos de rendimiento, no de diseño o características descriptivas»,
el borrador especificaba las \emph{funciones matemáticas} a computar,
sin constreñir cómo se computan. ISO estandarizó Classic McEliece en
junio de 2026, junto a ML-KEM y FrodoKEM, como enmienda a
ISO/IEC 18033-2 \cite{cmiso}. Mientras tanto, el NIST había
\emph{aplazado} la estandarización de McEliece en 2025, alegando que
«la estandarización concurrente por parte del NIST y de ISO crea el
riesgo de estándares incompatibles» y que «después de que concluya el
proceso de ISO, el NIST podría considerar desarrollar un estándar
basado en el de ISO» (traducción propia, \cite{cmnist}). El NIST, por
tanto, no ha estandarizado McEliece; su única selección de cifrado
post-cuántico adicional ha sido HQC \cite{nist8545}.
\end{historia}

\section{La variante \texttt{f}: forma semi-sistemática}\label{sec:param-f}

Recordemos cómo genera el capítulo~\ref{ch:keygen} la clave pública
(sección~«MatGen»): se reduce la matriz binaria $\hat H$ a forma
sistemática $(\mathrm{I}_{mt}\,|\,T)$ y, si la reducción falla, se
empieza de nuevo con otra clave. Con $(\mu,\nu)=(0,0)$ el intento
tiene éxito con probabilidad $\approx29\%$, así que en promedio hacen
falta unos $3{,}4$ intentos \cite{cmrationale}.

La forma \emph{$(\mu,\nu)$-semi-sistemática} relaja ese requisito. Una
matriz $R$ de rango $r$ en forma escalonada reducida, con pivotes en
las columnas $c_0<\cdots<c_{r-1}$, está en forma
$(\mu,\nu)$-semi-sistemática si no tiene filas nulas, $c_i=i$ para
$0\le i<r-\mu$, y cada pivote posterior admite un margen:
\[
  c_i\;\le\;i-\mu+\nu \qquad (0\le i<r).
\]
Para $(\mu,\nu)=(0,0)$ esto es, exactamente, la forma sistemática. Si
$\nu>\mu$, hay muchas más matrices admisibles, y con
$(\mu,\nu)=(32,64)$ basta con que $R$ sea escalonada salvo por los
últimos $\mu=32$ pivotes, que pueden desplazarse dentro de una ventana
de $\nu=64$ columnas. MatGen (caso general) reduce $\hat H$ hasta esa
forma, y luego intercambia las últimas columnas $i\leftrightarrow c_i$
---arrastrando los $\alpha_j$--- hasta dejar $(\mathrm{I}_{mt}\,|\,T)$.
El «libro de permutaciones» $c=(c_{mt-\mu},\dots,c_{mt-1})$ se guarda
en la clave privada (8 bytes; es la lista $c$ de la
tabla~\ref{tab:keygen-tamanos}, presente también en los conjuntos
$(\mu,\nu)=(0,0)$, donde codifica $2^{32}-1$).
La figura~\ref{fig:param-semisys} resume el mecanismo.

\begin{figure}[htbp]
\centering
\begin{tikzpicture}[
    cell/.style={draw=mcgray!50,minimum size=0.30cm,inner sep=0pt},
    one/.style={cell,fill=mcblue,draw=mcblue},
    piv/.style={cell,fill=mcorange,draw=mcorange},
    ilabel/.style={font=\scriptsize,align=center},
    arr/.style={-{Stealth[length=2.2mm]},thick,mcgray}
  ]
  % ventana sombreada
  \fill[mcorange2!30,rounded corners=2pt] (0.45,0.15) rectangle (2.25,-1.65);
  % rejilla izquierda 6x20
  \foreach \i in {0,...,5} \foreach \j in {0,...,19} {
    \node[cell] at (\j*0.3,-\i*0.3) {};
  }
  \foreach \i/\j in {0/0,1/1} \node[one] at (\j*0.3,-\i*0.3) {};
  \foreach \i/\j in {2/4,3/5,4/6,5/7} \node[piv] at (\j*0.3,-\i*0.3) {};
  \node[ilabel] at (1.35,0.55) {escalonada:};
  \node[ilabel,color=mcorange,align=center] at (4.6,-2.35)
    {ventana de $\nu$ columnas:\\ últimas $\mu$ filas};
  \draw[arr,mcorange] (3.3,-1.9) -- (2.05,-1.0);
  % flecha
  \draw[arr] (5.95,-0.75) -- (7.25,-0.75);
  \node[ilabel] at (6.6,-0.28) {$\le\mu$ intercambios};
  % rejilla derecha 6x12 (I|T)
  \fill[mcsky!45,rounded corners=2pt] (9.35,0.15) rectangle (11.05,-1.65);
  \foreach \i in {0,...,5} \foreach \j in {0,...,11} {
    \node[cell] at (7.6+\j*0.3,-\i*0.3) {};
  }
  \foreach \i in {0,...,5} \node[one] at (7.6+\i*0.3,-\i*0.3) {};
  \node[ilabel] at (9.3,0.55) {sistemática};
  \node[ilabel] at (9.3,-1.95) {$(\mathrm{I}_{mt}\,|\,T)$: clave pública};
  \node[ilabel,color=mcsky] at (11.5,-0.75) {$T$};
\end{tikzpicture}
\caption{La forma $(\mu,\nu)$-semi-sistemática, en esquema reducido
($\mu=4$, $\nu=6$; en el estándar, $(32,64)$). Izquierda: la matriz
escalonada; los primeros $mt-\mu$ pivotes ya están en la diagonal y
los últimos $\mu$ viven dentro de una ventana de $\nu$ columnas.
Basta con $\le\mu$ intercambios de columnas para alcanzar la forma
sistemática de la derecha.}
\label{fig:param-semisys}
\end{figure}

¿Qué se gana? Probabilidad de éxito. Con $(\mu,\nu)=(32,64)$ la
probabilidad de fallo de cada intento cae por debajo de $2^{-30}$: casi
siempre basta \emph{un} intento, frente a los $\approx3{,}4$ de
$(\mu,\nu)=(0,0)$. El precio: la computación en forma escalonada debe
hacerse sin ramificaciones dependientes de secretos (en tiempo
constante), y eso solo es razonable si queda confinada a una ventana
pequeña ---de ahí que $\nu$ se mantenga modesto---. La ganancia de
holgura en la demostración de seguridad también se cuantifica: aceptar
más matrices pierde como mucho $2$ bits de tensión en la reducción
\cite{cmrationale}.

¿Por qué existen entonces conjuntos sin \texttt{f}? Porque hay tres
argumentos para mantener $(\mu,\nu)=(0,0)$ \cite{cmrationale}: (1)~el
software optimizado de Classic McEliece ya acelera los reintentos
saltándose casi toda la eliminación gaussiana en los casos de fallo;
(2)~muchas aplicaciones no tienen la generación de claves como cuello
de botella; (3)~la simplicidad tiene valor. ISO estandarizó las
cuatro variantes para cada tamaño: quien quiera velocidad de
generación elige \texttt{f}; quien prefiera el camino más simple y
auditado, el conjunto base.

\section{La variante \texttt{pc}: confirmación de texto claro}\label{sec:param-pc}

En la ronda 3, la transformación CCA de Classic McEliece usaba
\emph{dos} defensas a la vez: rechazo implícito \emph{y} confirmación
de texto claro. La confirmación añade un segundo componente al
criptograma, $C_1=H(2,e)$, y el descifrado lo comprueba
\cite{cmpc}:

\begin{itemize}[leftmargin=1.5em,itemsep=3pt]
  \item $\textsf{Encap}$: $C_0=\textsf{Encode}(e,T)$,
        $C_1=H(2,e)$; el criptograma es $C=(C_0,C_1)$ y la clave,
        $K=H(1,e,C)$.
  \item $\textsf{Decap}$: se parte $C=(C_0,C_1)$; se descifra
        $C_0$ como siempre; después se recalcula $C_1'=H(2,e)$ y, si
        no coincide con $C_1$, también se toma la rama de rechazo
        ($e\leftarrow s$, $b\leftarrow0$). La clave sale de
        $K=H(b,e,C)$.
\end{itemize}

Nótese el detalle: la confirmación \emph{no} introduce errores
explícitos. Si la comprobación falla, \textsf{Decap} sigue devolviendo
$32$ bytes por la rama de rechazo implícito; el KEM continúa siendo
una función total. Lo que añade es una \emph{demostración}: $C_1$ solo
puede calcularlo quien conoce $e$, así que el criptograma certifica
que su emisor conocía el error. Quien posee la clave privada puede
verificar esa demostración explícitamente (descifra $C_0$, recalcula
$H(2,e')$ y compara con $C_1$), si su aplicación lo necesita.

El costo y la historia: $C_1$ ocupa $\lceil\ell/8\rceil=32$ bytes (los
mismos que una clave de sesión), y el hash pasa a tener tres tipos de
entrada, distinguidas por el primer byte: $(2,v)$, $(1,v,C)$ y
$(0,v,C)$. En la ronda 4, el equipo \emph{quitó} la confirmación del
conjunto por defecto y la dejó como variante opcional
\cite{cmmods3}, por tres razones: (1)~separarse limpiamente de la
patente estadounidense US~9.912.479 ---que no se aplica literalmente
a Classic McEliece, cuyo paquete de cifrado tiene un solo objeto, el
error, y no «el vector de error y el valor del texto claro», pero cuya
sombra legal se prefirió despejar, volviendo a la «forma tradicional
de Niederreiter»---; (2)~ahorrar los $32$ bytes, que en un esquema de
criptogramas diminutos se notan; y (3)~simplificar especificación y
software, ya que para las demostraciones ajustadas de IND-CCA2 en ROM
y cuasi-ajustadas en QROM (teoremas~\ref{thm:kem-rom}
y~\ref{thm:kem-qrom}) el rechazo implícito por sí solo basta.
En el debe: quitar la confirmación abre un ataque eficiente en un
modelo \emph{más allá} de IND-CCA2 en el que un solo bit de la clave
privada almacenada se invierte (por ejemplo, un fallo de DRAM); el
ataque análogo contra NTRU-HRSS \cite{ntrufault} motivó precisamente
este cuidado. La mitigación recomendada no toca el KEM: codificar la
clave privada almacenada con corrección de errores (los sistemas
comerciales ya hacen SECDED sobre la DRAM). ISO, en fin, estandariza
las dos familias de variantes: quien quiera las defensas duales de la
ronda 3 usa \texttt{pc} o \texttt{pcf}.

\section{Simply Decoded y Narrowly Decoded: el relleno}
\label{sec:param-narrowly}

Los conjuntos de parámetros no eligen un múltiplo de $8$ por
casualidad: el criptograma ocupa $\lceil mt/8\rceil$ bytes y la clave
pública $mt\lceil k/8\rceil$ bytes (una fila de $T$ por línea, rellena
cada una a byte completo). Cuando $mt$ o $k$ \emph{no} son múltiplos
de $8$, la última posición contiene bits de relleno que la
especificación deja a cero, y ahí nacen dos formas de leer una cadena
de bytes \cite{cmspec}:

\begin{itemize}[leftmargin=1.5em,itemsep=3pt]
  \item \textbf{Simply Decoded}: los bits de relleno se ignoran al
        interpretar la entrada;
  \item \textbf{Narrowly Decoded}: se rechazan las entradas
        (criptogramas y claves públicas) con relleno distinto de
        cero, devolviendo $\bot$.
\end{itemize}

De los dieciséis conjuntos, \textbf{solo \texttt{mceliece6960119}
tiene relleno}: allí $mt=1547$ (el criptograma tiene $194$ bytes, con
$5$ bits de relleno en el último) y $k=5413$ (cada fila de la clave
pública tiene $3$ bits de relleno). En los demás, $mt$ y $k$ son
múltiplos de $8$ y las dos lecturas coinciden. La
figura~\ref{fig:param-narrowly} ilustra el caso.

\begin{figure}[htbp]
\centering
\begin{tikzpicture}[
    bcell/.style={draw=mcgray!60,minimum size=0.5cm,inner sep=0pt,font=\scriptsize},
    rcell/.style={bcell,fill=mcblue!25},
    pcell/.style={bcell,fill=mcorange2!55},
    arr/.style={-{Stealth[length=2.2mm]},thick,mcgray},
    lbl/.style={font=\scriptsize,align=center}
  ]
  % bytes completos
  \node[draw=mcgray!60,rounded corners=2pt,minimum width=3.4cm,minimum height=0.5cm,
        font=\scriptsize,fill=mcgray!10] at (1.7,0) {bytes $0$--$192$};
  % ultimo byte
  \foreach \b in {0,...,7} {
    \pgfmathsetmacro{\cx}{4.15+\b*0.5}
    \ifnum\b<3 \node[rcell] at (\cx,0) {\b};
    \else \node[pcell] at (\cx,0) {\b};\fi
  }
  \node[lbl] at (5.9,0.55) {byte $193$ (el último): bits $1544$--$1551$};
  \node[rcell,minimum size=0.22cm] at (4.35,-0.75) {};
  \node[lbl,anchor=west] at (4.5,-0.75) {bits reales ($1547$ en total)};
  \node[pcell,minimum size=0.22cm] at (7.3,-0.75) {};
  \node[lbl,anchor=west] at (7.45,-0.75) {relleno ($5$ bits)};
  % interpretaciones
  \node[draw=mcgreen,fill=mcgreen2!45,rounded corners=2pt,font=\scriptsize,
        align=center,inner sep=4pt,text width=4.6cm] (sim) at (2.2,-2.5)
        {\textbf{Simply}: se ignoran los bits de relleno: el criptograma
         es el mismo elemento de $\Ftwo^{1547}$ y la clave no cambia.};
  \node[draw=mcred,fill=mcred2!45,rounded corners=2pt,font=\scriptsize,
        align=center,inner sep=4pt,text width=4.6cm] (nar) at (8.2,-2.5)
        {\textbf{Narrowly}: si algún bit de relleno es $1$, la entrada
         se rechaza ($\bot$).};
  \draw[arr] (5.4,-1.15) -- (3.6,-1.85);
  \draw[arr] (6.4,-1.15) -- (8.2,-1.85);
\end{tikzpicture}
\caption{El relleno de \texttt{mceliece6960119}: el criptograma son
$1547$ bits dentro de $194$ bytes, y a los $5$ bits finales la
especificación les da dos lecturas. Una implementación u otra debe
elegirse de forma consistente entre las dos partes del protocolo;
mezclarlas convierte en divergencia lo que es una maleabilidad
benigna.}
\label{fig:param-narrowly}
\end{figure}

¿Por qué se llama maleabilidad benigna? Porque dos cadenas de bytes
que difieran solo en el relleno representan el \emph{mismo}
criptograma (el mismo elemento de $\Ftwo^{mt}$), y por tanto el mismo
$K$ en modo Simply; en modo Narrowly una se acepta y la otra no. La
distinción no rompe la seguridad ---el teorema de IND-CCA2 cubre las
cadenas codificadas---, pero es un contrato de interoperabilidad que
las partes deben acordar.

\begin{trampa}[Tres bits por fila, $580$ bytes]
Un implementador que empaquete la clave pública «apretada», como una
tira continua de $kmt$ bits, obtendrá para \texttt{mceliece6960119}
una clave de $\lceil kmt/8\rceil=1\,046\,739$ bytes. La cadena
estándar tiene $1\,047\,319$: la diferencia, $580$ bytes, son los
$3$ bits de relleno de cada una de las $1547$ filas
($1547\times3/8=580{,}125$). Ambos empaquetados describen la misma
matriz $T$, pero solo el segundo es válido frente a otra
implementación ---y frente a los vectores de prueba del
estándar---. Con $T$ en la mano el error se detecta enseguida, pero es
un clásico de los primeros días de una implementación.
\end{trampa}

\section{¿Cuánta seguridad?}\label{sec:param-seguridad}

La tabla~\ref{tab:param-seguridad} resume las estimaciones de la guía
de seguridad \cite{cmsecurity}: el mejor ataque de decodificación por
conjuntos de información (capítulo~\ref{ch:isd}) según el estimador de
Esser--Bellini, en dos de sus modelos de memoria, junto a la categoría
que cada conjunto reclama.

\begin{table}[htbp]
\centering
\footnotesize
\setlength{\tabcolsep}{5pt}
\begin{tabular}{lccccc}
\toprule
conjunto & Prange & MO (mem.\ libre) & MO (mem.\ con costo) & categoría & categoría\\
 & & & & (mem.\ realista) & (mem.\ gratis)\\
\midrule
\texttt{mceliece348864} & $173{,}4$ & $140{,}8$ & $158{,}6$ & 1 & 1\\
\texttt{mceliece460896} & $216{,}8$ & $179{,}8$ & $201{,}0$ & 3 & 2\\
\texttt{mceliece6688128} & $295{,}7$ & $246{,}2$ & $278{,}2$ & 5 & 4\\
\texttt{mceliece6960119} & $296{,}8$ & $245{,}7$ & $279{,}2$ & 5 & 4\\
\texttt{mceliece8192128} & $334{,}1$ & $275{,}6$ & $315{,}6$ & 5 & 5\\
\bottomrule
\end{tabular}
\caption{Costos en bits ($\log_2$) de los mejores ataques conocidos
según el estimador de Esser--Bellini: contra Prange (1962) y contra
May--Ozerov (MO), con memoria de acceso gratis y con costo de acceso
en raíz cuadrada. Las dos últimas columnas: la categoría que el
equipo asigna contando costos realistas de memoria, y la que
correspondería si la memoria se contara como gratis.}
\label{tab:param-seguridad}
\end{table}

Tres advertencias sobre esa tabla, todas de la propia guía:

\begin{enumerate}[leftmargin=1.6em,itemsep=3pt]
  \item \textbf{Los números subestiman.} El estimador cuenta cada
        operación vectorial como \emph{una} operación, sea cual sea su
        longitud. El $140{,}8$ de \texttt{mceliece348864}, por ejemplo,
        cuenta más de $2^{143}$ operaciones de bits, y el ataque
        requiere memoria $2^{86{,}6}$ \cite{cmsecurity}.
  \item \textbf{La memoria no es gratis.} La columna «mem.\ libre»
        supone acceso libre a memoria arbitrariamente grande,
        físicamente irreal; ahí los algoritmos con
        «representaciones» (BJMM, MO) brillan. Cuando se cobra la
        memoria a su costo plausible, la ventaja se estrecha.
  \item \textbf{Las categorías dependen del modelo.} Contando
        memoria realista, el equipo asigna $1,3,5,5,5$; contando
        memoria gratis y operaciones de bits, $1,2,4,4,5$. Ambas
        tablas son sostenibles sobre los mismos datos; el número clave
        es que \texttt{mceliece6960119} es categoría 5 si se cuentan
        los accesos a memoria ($2^{279{,}2}>2^{272}$, la referencia de
        AES-256) y categoría 4 si no ($2^{245{,}7}<2^{256}$).
\end{enumerate}

Las referencias del NIST para las categorías 1, 3 y 5 son búsquedas
exhaustivas de AES-128, AES-192 y AES-256, estimadas en $2^{143}$,
$2^{207}$ y $2^{272}$ operaciones de bits. De ahí las tensiones
numéricas reseñadas: $140{,}8$, el peor valor, roza por abajo la
referencia de la categoría 1; $201{,}0$ cae bajo $2^{207}$; y
$245{,}7$, bajo $2^{256}$. El equipo responde, con razón documentada,
que estas comparaciones mezclan métricas: no es lo mismo «operaciones
de bits contadas por un estimador que trata un XOR de mil bytes como
una operación» que operaciones reales. La discusión completa, y por
qué el récord de estabilidad del esquema es el argumento más fuerte,
llega en los capítulos~\ref{ch:seguridad} y~\ref{ch:isd}.

\section{Elegir un conjunto}\label{sec:param-elegir}

Con todo lo anterior, la guía práctica es corta:

\begin{itemize}[leftmargin=1.5em,itemsep=3pt]
  \item \textbf{Tamaño.} Para seguridad a largo plazo, el equipo
        recomienda los tamaños \texttt{mceliece6*}
        ($6688128$, $6960119$) \cite{cmiso}; añaden un margen
        generoso sobre la categoría~5 clásica. Si el entorno aprieta
        (clave pública de $512$~KiB), el candidato es
        \texttt{mceliece460896}; \texttt{mceliece348864} existe en la
        especificación, pero queda fuera de ISO y no es la elección
        para décadas.
  \item \textbf{Variante \texttt{f}.} Úsese cuando la generación de
        claves sea cuello de botella o deba ser rápida en tiempo
        constante; añade la forma semi-sistemática sin cambiar
        ningún tamaño.
  \item \textbf{Variante \texttt{pc}/\texttt{pcf}.} Úsese cuando el
        protocolo quiera la confirmación de texto claro de la ronda 3
        (por ejemplo, para verificar el conocimiento de $e$ o para
        defensas en modelos con fallos de clave privada); cuesta
        $32$ bytes por criptograma.
\end{itemize}

Para calibrar los compromisos, la tabla~\ref{tab:param-contexto}
compara los tamaños con los otros KEM que un despliegue real barajará.
La silueta de Classic McEliece es inconfundible: una clave pública
doscientas veces mayor que la de ML-KEM, un criptograma entre cinco y
ocho veces menor, y una clave privada moderada. La clave de HQC, en
cambio, es de $40$ bytes: es una semilla, regenerada al vuelo, la
misma filosofía que la $\delta$ de Classic McEliece pero llevada al
secreto entero.

\begin{table}[htbp]
\centering
\small
\setlength{\tabcolsep}{7pt}
\begin{tabular}{lrrr}
\toprule
esquema & $ek$ (bytes) & $ct$ (bytes) & $sk$ (bytes)\\
\midrule
ML-KEM-768 & 1184 & 1088 & 2400\\
ML-KEM-1024 & 1568 & 1568 & 3168\\
HQC-128 & 2249 & 4497 & 40\\
HQC-256 & 7245 & 14485 & 40\\
\texttt{mceliece460896} & \num{524160} & 156 & \num{13608}\\
\texttt{mceliece6960119} & \num{1047319} & 194 & \num{13948}\\
\bottomrule
\end{tabular}
\caption{Tamaños de los KEM post-cuánticos relevantes (ML-KEM según
FIPS~203 \cite{fips203}; HQC según el informe del NIST \cite{nist8545};
Classic McEliece según la tabla~\ref{tab:param-base}). El criptograma
de McEliece es el más pequeño con diferencia; la clave pública, la más
grande.}
\label{tab:param-contexto}
\end{table}

\begin{practica}[¿Y en un protocolo real?]
La clave pública de Classic McEliece no viaja en cada conexión: viaja
una vez y se cachea (es estable durante años). El criptograma, en
cambio, viaja en cada intercambio: ahí el tamaño importa, y es donde
los $96$--$240$ bytes juegan a favor. En TLS~1.3, por ejemplo, ML-KEM
y McEliece pueden convivir como grupos híbridos; la página de
implementación del equipo documenta los detalles
\cite{cmimpl}.
\end{practica}

\section{Ejercicios}\label{sec:param-ejercicios}

\begin{enumerate}[leftmargin=1.6em,itemsep=4pt]
  \item \textbf{Aritmética de bytes.} Con $ct=\lceil mt/8\rceil$ y
        $ek=mt\lceil k/8\rceil$, verifique las columnas de la
        tabla~\ref{tab:param-base}. ¿Por qué para cuatro conjuntos
        $ek=\lceil kmt/8\rceil$ pero para \texttt{mceliece6960119} no?
        Calcule la diferencia exacta.
  \item \textbf{La variante \texttt{pc}.} Para
        \texttt{mceliece460896pc}, complete la
        tabla~\ref{tab:param-iso} para los dieciséis conjuntos y
        calcule después cuántos bytes recibe cada hash: $H(2,e)$ por
        un lado, y $H(1,e,C)$ y $H(0,s,C)$ por otro.
  \item \textbf{El listón cuántico de ISO.} Usando el modelo de raíz
        cuadrada, calcule la mitad del logaritmo del ataque de Stern
        (tabla~\ref{tab:param-seguridad}) para los cinco conjuntos y
        compruebe qué conjuntos superan $2^{128}$. ¿Coincide con los
        que eligió el borrador ISO? ¿Y con lo que ISO decidió
        estandarizar al final? Explique la diferencia.
  \item \textbf{Intentos de generación.} Con $(\mu,\nu)=(0,0)$ el
        éxito de MatGen ronda el $29\%$: calcule el número esperado de
        intentos y la probabilidad de necesitar más de dos. Con
        $(\mu,\nu)=(32,64)$ el fallo por intento es $<2^{-30}$:
        calcule el número esperado de intentos. ¿Cuánto trabajo de
        generación se ahorra la variante \texttt{f}?
  \item \textbf{Relleno benigno.} Para \texttt{mceliece6960119},
        construya dos cadenas de $194$ bytes que difieran solo en los
        $5$ bits de relleno. ¿Cuántos $K$ distintos producen en modo
        Simply? ¿Y en modo Narrowly? ¿Por qué esta diferencia debe
        resolverse por acuerdo de protocolo y no dejarse al azar de
        la implementación?
  \item \textbf{Elegir con criterio.} Un servicio de mensajería
        cachea claves públicas de servidor y envía millones de
        criptogramas pequeños; una etiqueta IoT de bajo consumo
        genera su propia clave una vez y no la cachea. ¿Qué conjunto
        y qué variante propondría para cada caso? Justifique con las
        tablas~\ref{tab:param-base} y~\ref{tab:param-contexto}.
\end{enumerate}
